\documentclass[../../main.tex]{subfiles}
\begin{document}

\begin{center}
  {\fontsize{48pt}{58pt}\selectfont \textbf{2005}\normalsize\textbf{年度}}
\end{center}

\begin{table}[htb]
  \centering
  \caption{出題テーマと難易度}
  \begin{tabular}{|c|p{9cm}|c|}
    \hline
    \textbf{問題番号} & \textbf{テーマ} & \textbf{難易度} \\
    \hline
    第1問 & 放物線と線分が交わる係数の範囲 & \\
    \hline
    第2問 & 対数を使う自然数の個数 & \\
    \hline
    第3問 & 複素数の和と平方和で定まる三角形 & \\
    \hline
    第4問 & 整数の立方差 & \\
    \hline
    第5問 & 余弦曲線と有理関数の共有点 & \\
    \hline
    第6問 & 隣接車両の色の塗り方 & \\
    \hline
  \end{tabular}
\end{table}

\begin{table}[htb]
  \centering
  \caption{解答}
  \renewcommand{\arraystretch}{1.4}
  \begin{tabular}{|c|c|p{9cm}|}
    \hline
    \textbf{問題番号} & \textbf{小問} & \textbf{答え} \\
    \hline
    第1問 & -- & $b(a+b-1)\le0$，または$0\le a\le2$かつ$b\ge0$，$a+b\ge1$，$b\le(a-2)^2/4$ \\
    \hline
    第2問 & -- & $31$個（$n=32,\ldots,62$） \\
    \hline
    第3問 & -- & 重心が原点の正三角形 \\
    \hline
    第4問 & -- & $(a,b)=(9,8),(-8,-9),(6,-1),(1,-6)$ \\
    \hline
    第5問 & (1) & 共有点の接線は$(0,1)$を通る（証明問題） \\
    \cline{2-3}
    & (2) & 共有点はただ1つ（証明問題） \\
    \hline
    第6問 & -- & $\dfrac{2^{n+2}+(-1)^{n-1}}{3}$通り \\
    \hline
  \end{tabular}
\end{table}

\paragraph{試験形態}
6問．試験時間は要確認．

\end{document}
